\paragraph{审计偶数窗上的最小退化扇区、群胚维数、规范群熵与黄金密度分裂}
在定理 \ref{thm:xi-foldbin-minimal-center-valued-index-spectrum-top} 已把纤维多重度场压缩为中心谱变量之后，审计偶数窗
\[
m\in\{6,8,10,12\}
\]
上由最小退化扇区驱动的碰撞、Abel 化、群胚中心与规范群熵还可进一步收束为一组完全对象层结论。沿用定义 \ref{def:fold-bin-degeneracy-sectors}、定理 \ref{thm:fold-bin-min-degeneracy-fib-audited}、推论 \ref{cor:fold-bin-minsector-complement-maxfiber-audited}、定理 \ref{thm:xi-foldbin-collision-probability-normalized-groupoid-dimension}、定理 \ref{thm:xi-foldbin-parity-charge-split-exact-minrank}、命题 \ref{prop:fold-bin-groupoid-algebra-wedderburn}、注记 \ref{rem:fold-bin-gauge-entropy-scale} 与推论 \ref{cor:fold-bin-minsector-density-phi-minus2} 的记号，记
\[
d_*:=d_{\min}(m)=F_{m/2},
\qquad
\mathcal S_{m,*}:=\mathcal S_{m,d_*},
\qquad
r_m:=|X_m\setminus \mathcal S_{m,*}|=D_{2m-2}=F_{m+1}.
\]
为避免与定理 \ref{thm:xi-foldbin-minimal-center-valued-index-spectrum-top} 中的谱顶端记号 \(M_m\) 混淆，记补扇区总质量
\[
T_m:=2^m-F_m d_*,
\qquad
T_m=q_m r_m+s_m,\qquad 0\le s_m<r_m.
\]
并置
\[
A_m:=\CC[\mathcal R_m^{\mathrm{bin}}],
\qquad
A_{m,\min}:=\bigoplus_{w\in\mathcal S_{m,*}}M_{d_*}(\CC),
\qquad
A_{m,\mathrm{out}}:=\bigoplus_{w\in X_m\setminus\mathcal S_{m,*}}
M_{d_m^{\mathrm{bin}}(w)}(\CC).
\]

\begin{theorem}[群胚代数维数与碰撞矩的严格等同及最小扇区平衡下界]\label{thm:xi-audited-even-groupoid-dimension-collision-balanced-lower-bound}
有严格恒等式
\[
\boxed{\;
\dim_{\CC}A_m
=
\sum_{w\in X_m}\bigl(d_m^{\mathrm{bin}}(w)\bigr)^2
=
4^m\,\mathsf{Col}^{\mathrm{bin}}_m.\;}
\]
进一步，
\[
\boxed{\;
\dim_{\CC}A_m
\ge
F_m d_*^2+(r_m-s_m)q_m^2+s_m(q_m+1)^2.\;}
\]
并且等号成立当且仅当补扇区中的全部纤维大小恰只取 \(q_m\) 与 \(q_m+1\) 两值，其中恰有 \(s_m\) 条取 \(q_m+1\)。
\end{theorem}
\begin{proof}
由定理 \ref{thm:xi-foldbin-collision-probability-normalized-groupoid-dimension}，
\[
\dim_{\CC}A_m
=
\sum_{w\in X_m}\bigl(d_m^{\mathrm{bin}}(w)\bigr)^2
=
4^m\,\mathsf{Col}^{\mathrm{bin}}_m,
\]
第一条结论成立。

再证下界。由定理 \ref{thm:fold-bin-min-degeneracy-fib-audited}，
\[
|\mathcal S_{m,*}|=F_m,
\qquad
d_m^{\mathrm{bin}}(w)=d_*
\quad (w\in\mathcal S_{m,*}),
\]
而由推论 \ref{cor:fold-bin-minsector-complement-maxfiber-audited}，
\[
|X_m\setminus \mathcal S_{m,*}|=r_m.
\]
将补扇区上的纤维大小重记为
\[
e_1,\dots,e_{r_m},
\]
则
\[
e_i\ge d_*+1,
\qquad
\sum_{i=1}^{r_m}e_i=T_m,
\]
并且
\[
\dim_{\CC}A_m
=
F_m d_*^2+\sum_{i=1}^{r_m}e_i^2.
\]
因此只需在固定长度 \(r_m\) 与固定总和 \(T_m\) 的约束下最小化
\(
\sum_i e_i^2
\)。

若某一对 \(a,b\) 满足 \(a\ge b+2\)，则
\[
a^2+b^2-(a-1)^2-(b+1)^2
=
2(a-b)-2
>0.
\]
故在保持总和不变的前提下，用 \((a-1,b+1)\) 代替 \((a,b)\) 会严格降低平方和。不断迭代可知，极小值只能在所有坐标彼此至多相差 \(1\) 时取得。把
\[
T_m=q_m r_m+s_m,\qquad 0\le s_m<r_m
\]
代入，唯一可能的极小配置即为
\[
\underbrace{q_m,\dots,q_m}_{r_m-s_m\text{ 次}},
\underbrace{q_m+1,\dots,q_m+1}_{s_m\text{ 次}}.
\]
于是
\[
\sum_{i=1}^{r_m}e_i^2
\ge
(r_m-s_m)q_m^2+s_m(q_m+1)^2,
\]
从而得到所示下界。

最后说明等号条件。对四个审计偶数窗直接代入可得
\[
q_m\in\{3,5,8,12\},
\qquad
d_*\in\{2,3,5,8\},
\]
故总有 \(q_m\ge d_*+1\)，从而上述平衡配置确为补扇区约束 \(e_i\ge d_*+1\) 下的可行点。于是等号成立当且仅当补扇区中的全部纤维大小恰为 \(q_m\) 与 \(q_m+1\) 两值，且其中有 \(s_m\) 条取 \(q_m+1\)。证毕。
\end{proof}

\begin{theorem}[审计偶数窗上的 Abel 化与群胚代数中心的 Fibonacci 双尺度分裂]\label{thm:xi-audited-even-abel-center-fibonacci-splitting}
有规范直和分裂
\[
\boxed{\;
\bigl(G_m^{\mathrm{bin}}\bigr)^{\mathrm{ab}}
\cong
(\ZZ/2\ZZ)^{\mathcal S_{m,*}}
\oplus
(\ZZ/2\ZZ)^{X_m\setminus\mathcal S_{m,*}}
\cong
(\ZZ/2\ZZ)^{F_m}\oplus(\ZZ/2\ZZ)^{D_{2m-2}}.\;}
\]
并且
\[
\boxed{\;
Z(A_m)
\cong
Z(A_{m,\min})\oplus Z(A_{m,\mathrm{out}})
\cong
\CC^{F_m}\oplus\CC^{D_{2m-2}}.\;}
\]
其中第一直和因子恰由最小退化扇区 \(\mathcal S_{m,*}\) 提供，第二直和因子恰由其补集提供。
\end{theorem}
\begin{proof}
由定理 \ref{thm:xi-foldbin-parity-charge-split-exact-minrank}，
\[
\bigl(G_m^{\mathrm{bin}}\bigr)^{\mathrm{ab}}
\cong
(\ZZ/2\ZZ)^{N_m},
\qquad
N_m:=\#\{w\in X_m:\ d_m^{\mathrm{bin}}(w)\ge 2\}.
\]
而在审计偶数窗上，定理 \ref{thm:fold-bin-min-degeneracy-fib-audited} 给出
\[
d_{\min}(m)=d_*\ge 2,
\]
故每条纤维都满足 \(d_m^{\mathrm{bin}}(w)\ge 2\)，从而
\[
N_m=|X_m|.
\]
于是
\[
\bigl(G_m^{\mathrm{bin}}\bigr)^{\mathrm{ab}}
\cong
(\ZZ/2\ZZ)^{X_m}.
\]
再由集合分拆
\[
X_m=\mathcal S_{m,*}\sqcup(X_m\setminus\mathcal S_{m,*}),
\]
得到
\[
(\ZZ/2\ZZ)^{X_m}
\cong
(\ZZ/2\ZZ)^{\mathcal S_{m,*}}
\oplus
(\ZZ/2\ZZ)^{X_m\setminus\mathcal S_{m,*}}.
\]
最后由
\[
|\mathcal S_{m,*}|=F_m,
\qquad
|X_m\setminus\mathcal S_{m,*}|=D_{2m-2},
\]
便得第一条分裂。

对中心，由命题 \ref{prop:fold-bin-groupoid-algebra-wedderburn}，
\[
A_m
\cong
\bigoplus_{w\in X_m}M_{d_m^{\mathrm{bin}}(w)}(\CC)
=
A_{m,\min}\oplus A_{m,\mathrm{out}}.
\]
有限维半单代数的中心对直和可加，因此
\[
Z(A_m)\cong Z(A_{m,\min})\oplus Z(A_{m,\mathrm{out}}).
\]
又每个矩阵块的中心均为 \(\CC\)，故
\[
Z(A_{m,\min})\cong \CC^{\mathcal S_{m,*}},
\qquad
Z(A_{m,\mathrm{out}})\cong \CC^{X_m\setminus\mathcal S_{m,*}}.
\]
取维数并代入上述两组基数，即得
\[
Z(A_m)\cong \CC^{F_m}\oplus \CC^{D_{2m-2}}.
\]
证毕。
\end{proof}

\begin{theorem}[规范群熵的最优平衡下界与补扇区熵缺口]\label{thm:xi-audited-even-gauge-entropy-balanced-lower-bound}
沿用注记 \ref{rem:fold-bin-gauge-entropy-scale} 的 bits 记号
\[
H_m^{\mathrm{gauge}}
:=
\log_2|G_m^{\mathrm{bin}}|
=
\sum_{w\in X_m}\log_2\bigl(d_m^{\mathrm{bin}}(w)!\bigr).
\]
则有最优下界
\[
\boxed{\;
H_m^{\mathrm{gauge}}
\ge
F_m\log_2(d_*!)
+
(r_m-s_m)\log_2(q_m!)
+
s_m\log_2((q_m+1)!).\;}
\]
并且等号成立当且仅当补扇区中的全部纤维大小恰只取 \(q_m\) 与 \(q_m+1\) 两值。

等价地，若定义补扇区熵缺口
\[
\Delta H_m:=H_m^{\mathrm{gauge}}-F_m\log_2(d_*!),
\]
则
\[
\boxed{\;
\Delta H_m
\ge
(r_m-s_m)\log_2(q_m!)
+
s_m\log_2((q_m+1)!).\;}
\]
\end{theorem}
\begin{proof}
由定义，最小扇区上的每条纤维大小都等于 \(d_*\)，其条数为 \(F_m\)。把补扇区上的纤维大小仍记为
\[
e_1,\dots,e_{r_m},
\qquad
\sum_{i=1}^{r_m}e_i=T_m.
\]
则
\[
H_m^{\mathrm{gauge}}
=
F_m\log_2(d_*!)
+
\sum_{i=1}^{r_m}\log_2(e_i!).
\]

定义
\[
f(n):=\log_2(n!)
\qquad (n\in\ZZ_{\ge 1}).
\]
则
\[
f(n+1)-f(n)=\log_2(n+1),
\]
故差分严格递增，亦即 \(f\) 在正整数上严格离散凸。若某一对 \(a,b\) 满足 \(a\ge b+2\)，则
\[
f(a)+f(b)-f(a-1)-f(b+1)
=
\log_2 a-\log_2(b+1)
>0.
\]
因此在固定总和 \(T_m\) 与固定项数 \(r_m\) 的约束下，\(\sum_i f(e_i)\) 的极小值同样由“尽可能均匀”的配置取得，也就是全部 \(e_i\) 只取
\[
q_m,\qquad q_m+1
\]
两值，其中恰有 \(s_m\) 个取 \(q_m+1\)。由定理 \ref{thm:xi-audited-even-groupoid-dimension-collision-balanced-lower-bound} 证明末尾的可行性说明，该配置在四个审计偶数窗上确为可行。故
\[
\sum_{i=1}^{r_m}\log_2(e_i!)
\ge
(r_m-s_m)\log_2(q_m!)
+
s_m\log_2((q_m+1)!),
\]
代回即得第一条下界。

若取等，则上述离散平衡过程中每一步都不能继续降低函数值，于是全部 \(e_i\) 必彼此至多相差 \(1\)，亦即只可能取 \(q_m\) 与 \(q_m+1\) 两值；反之若补扇区纤维正好取这两值，则显然达到下界。第二条关于 \(\Delta H_m\) 的公式只是把最小扇区项移至左端。证毕。
\end{proof}

\begin{theorem}[子覆盖密度延拓下的 Abel 秩密度与中心密度黄金分裂]\label{thm:xi-audited-even-minsector-algebraic-density-golden-splitting}
若“子覆盖同构”恒等式
\[
|\mathcal S_{m,*}|=|X_{m-2}|
\]
沿某个趋于无穷的偶数序列成立，则在同一序列上有
\[
\boxed{\;
\frac{
\operatorname{rank}_{\FF_2}\bigl((\ZZ/2\ZZ)^{\mathcal S_{m,*}}\bigr)
}{
\operatorname{rank}_{\FF_2}\!\bigl((G_m^{\mathrm{bin}})^{\mathrm{ab}}\bigr)
}
=
\frac{\dim_{\CC}Z(A_{m,\min})}{\dim_{\CC}Z(A_m)}
=
\frac{F_m}{F_{m+2}}
\longrightarrow
\varphi^{-2}.\;}
\]
并且补扇区的对应密度满足
\[
\boxed{\;
\frac{
\operatorname{rank}_{\FF_2}\bigl((\ZZ/2\ZZ)^{X_m\setminus\mathcal S_{m,*}}\bigr)
}{
\operatorname{rank}_{\FF_2}\!\bigl((G_m^{\mathrm{bin}})^{\mathrm{ab}}\bigr)
}
=
\frac{\dim_{\CC}Z(A_{m,\mathrm{out}})}{\dim_{\CC}Z(A_m)}
=
\frac{F_{m+1}}{F_{m+2}}
\longrightarrow
\varphi^{-1}.\;}
\]
\end{theorem}
\begin{proof}
由定理 \ref{thm:xi-audited-even-abel-center-fibonacci-splitting}，
\[
\operatorname{rank}_{\FF_2}\bigl((\ZZ/2\ZZ)^{\mathcal S_{m,*}}\bigr)
=
|\mathcal S_{m,*}|,
\qquad
\operatorname{rank}_{\FF_2}\!\bigl((G_m^{\mathrm{bin}})^{\mathrm{ab}}\bigr)
=
|X_m|,
\]
并且
\[
\dim_{\CC}Z(A_{m,\min})=|\mathcal S_{m,*}|,
\qquad
\dim_{\CC}Z(A_m)=|X_m|.
\]
因此最小扇区的 Abel 秩密度与中心密度完全一致，均等于
\[
\frac{|\mathcal S_{m,*}|}{|X_m|}.
\]
在假设下，推论 \ref{cor:fold-bin-minsector-density-phi-minus2} 已证明
\[
\frac{|\mathcal S_{m,*}|}{|X_m|}
\longrightarrow
\varphi^{-2},
\]
故第一条极限成立。

对补扇区，同理有
\[
\frac{
\operatorname{rank}_{\FF_2}\bigl((\ZZ/2\ZZ)^{X_m\setminus\mathcal S_{m,*}}\bigr)
}{
\operatorname{rank}_{\FF_2}\!\bigl((G_m^{\mathrm{bin}})^{\mathrm{ab}}\bigr)
}
=
\frac{|X_m\setminus\mathcal S_{m,*}|}{|X_m|}
=
\frac{F_{m+1}}{F_{m+2}},
\]
以及
\[
\frac{\dim_{\CC}Z(A_{m,\mathrm{out}})}{\dim_{\CC}Z(A_m)}
=
\frac{|X_m\setminus\mathcal S_{m,*}|}{|X_m|}
=
\frac{F_{m+1}}{F_{m+2}}.
\]
再由
\[
\frac{F_{m+1}}{F_{m+2}}
=
1-\frac{F_m}{F_{m+2}}
\longrightarrow
1-\varphi^{-2}
=
\varphi^{-1},
\]
即得第二条极限。证毕。
\end{proof}

由此，在审计偶数窗上，最小退化扇区不再只是一个局域计数对象，而是同时控制以下四层结构：
\[
\mathcal S_{m,*}
\Longrightarrow
\dim_{\CC}A_m=4^m\mathsf{Col}^{\mathrm{bin}}_m
\Longrightarrow
\bigl(G_m^{\mathrm{bin}}\bigr)^{\mathrm{ab}}\ \text{与}\ Z(A_m)\ \text{的 Fibonacci 双裂解}
\Longrightarrow
H_m^{\mathrm{gauge}}\ \text{的最优平衡下界}
\Longrightarrow
(\varphi^{-2},\varphi^{-1})\ \text{型代数密度分裂}.
\]

\endinput
